% chapters/02_mathematical_preliminaries.tex

\section{Overview}

This chapter develops the complete mathematical toolkit required for all subsequent
chapters. No prior knowledge of optimal transport or stochastic partial differential
equations is assumed; every result is stated with either a complete proof or a detailed
proof sketch, with references to the authoritative sources for full details. The
standing assumptions A1--A7 referenced throughout are collected in Section 1.9 of
Chapter 1.

The chapter covers six interlocking topics:

\begin{enumerate}
    \item \textbf{Wasserstein-2 space} $(\mathcal{P}_2(\mathbb{R}^d), W_2)$: optimal
    transport, metric space structure, geodesics, weak convergence, rates of
    convergence for empirical measures, and the Lions derivative.
    \item \textbf{Brownian motion and It\^o calculus}: construction, quadratic
    variation, the It\^o integral, It\^o's formula, and stochastic differential
    equations.
    \item \textbf{Fokker--Planck SPDE}: derivation from It\^o's formula, weak
    formulation, strong form, Sobolev regularity, and the Gelfand triple framework. We
    clarify the distinction between the linear Fokker--Planck SPDE and the nonlinear
    Kushner--Zakai equation.
    \item \textbf{Martingale theory}: martingales, stopping times, Doob's inequality,
    Burkholder--Davis--Gundy, martingale representation, and Girsanov's
    theorem.
    \item \textbf{Sobolev spaces and functional analysis}: Sobolev embedding,
    Hahn--Banach extension, Lumer--Phillips theorem, spectral gap, and log-Sobolev
    inequalities.
    \item \textbf{Stochastic Maximum Principle}: the control problem, Hamiltonian,
    adjoint BSDE, optimality conditions, and connections to Hamilton--Jacobi--Bellman
    equations.
\end{enumerate}

A summary table at the end of the chapter (\Cref{sec:mp-summary}) maps each tool to its
first use in the thesis, serving as a quick reference. For readers seeking a deeper
treatment, we recommend Villani (2009) for optimal transport, Carmona--Delarue (2018,
Vol.\ I, Chapters 1--2) for Wasserstein analysis in the MFG context, and
Karatzas--Shreve (1991) for stochastic calculus.


\section{Wasserstein-2 Space}
\label{sec:mp-wasserstein}

\subsection{Probability Measures with Finite Second Moment}

Let $\mathcal{B}(\mathbb{R}^d)$ denote the Borel $\sigma$-algebra on $\mathbb{R}^d$. The
space $\mathcal{P}_2(\mathbb{R}^d)$ consists of all Borel probability measures $\mu$ on
$\mathbb{R}^d$ with finite second moment:
\begin{equation}
    \mathcal{P}_2(\mathbb{R}^d) := \left\{ \mu \in \mathcal{P}(\mathbb{R}^d) \;\middle|\;
    \int_{\mathbb{R}^d} |x|^2\, \mu(dx) < \infty \right\}.
    \label{eq:p2-def}
\end{equation}

Elements of $\mathcal{P}_2(\mathbb{R}^d)$ arise naturally in our setting:
\begin{itemize}
    \item As the law $\mu = \mathcal{L}(X)$ of an $\mathbb{R}^d$-valued random variable
    with $\mathbb{E}[|X|^2] < \infty$.
    \item As the empirical measure $\mu^N = \frac{1}{N}\sum_{i=1}^N \delta_{x_i}$ when
    the particles satisfy $\frac{1}{N}\sum_i |x_i|^2 < \infty$.
    \item As the mean-field distribution $\mu_t$ in the MFG system, provided the
    agents' positions remain square-integrable in expectation.
\end{itemize}

\subsection{Optimal Transport and the Wasserstein Distance}

\begin{definition}[Coupling and Wasserstein-2 Distance]
A coupling of $\mu, \nu \in \mathcal{P}_2(\mathbb{R}^d)$ is a probability measure
$\pi \in \mathcal{P}(\mathbb{R}^d \times \mathbb{R}^d)$ with marginals $\mu$ and $\nu$:
$\pi(\cdot \times \mathbb{R}^d) = \mu$, $\pi(\mathbb{R}^d \times \cdot) = \nu$. The set of
all couplings is denoted $\Pi(\mu, \nu)$. The Wasserstein-2 distance between $\mu$ and
$\nu$ is
\begin{equation}
    W_2(\mu, \nu) := \left( \inf_{\pi \in \Pi(\mu,\nu)} \int_{\mathbb{R}^d \times
    \mathbb{R}^d} |x-y|^2\, \pi(dx,dy) \right)^{1/2}.
    \label{eq:wasserstein-def}
\end{equation}
\end{definition}

The infimum in \eqref{eq:wasserstein-def} is attained: there exists an optimal coupling
$\pi^* \in \Pi(\mu,\nu)$ achieving the infimum (Villani, 2009, Theorem 5.10). The
optimal $\pi^*$ is supported on the subdifferential of a convex function, the Brenier
potential.

\begin{theorem}[Brenier's Theorem, 1991]
\label{thm:brenier}
Let $\mu, \nu \in \mathcal{P}_2(\mathbb{R}^d)$ with $\mu$ absolutely continuous with
respect to Lebesgue measure. Then there exists a unique (up to $\mu$-a.e.\ equality)
convex function $\varphi : \mathbb{R}^d \to \mathbb{R}$ such that the map
$T = \nabla \varphi$ satisfies $T_\# \mu = \nu$ (i.e., $\nu(A) = \mu(T^{-1}(A))$ for all
$A$) and
\[
    W_2(\mu,\nu)^2 = \int_{\mathbb{R}^d} |x - \nabla\varphi(x)|^2\, \mu(dx).
\]
The optimal coupling is $\pi^* = (\mathrm{id} \times T)_\# \mu$, i.e., the law of
$(X, T(X))$ when $X \sim \mu$.
\end{theorem}

\begin{proof}[Proof sketch]
The proof proceeds in three steps: (i) Kantorovich duality reduces the problem to
finding a pair of conjugate convex functions $(\varphi, \varphi^*)$; (ii) the
optimality condition $y \in \partial\varphi(x)$ for $\pi^*$-a.e.\ $(x,y)$ implies the
transport is a gradient map; (iii) uniqueness follows from strict convexity of the cost
$c(x,y) = |x-y|^2$. See Villani (2009, Ch.\ 9) for the complete argument.
\end{proof}

\begin{remark}[Discrete measures and regularisation]
Brenier's Theorem requires absolute continuity of the source measure $\mu$. In
particle-method applications (Chapter 5), the empirical measure $\mu_t^N$ is a sum of
Dirac masses, and the equilibrium measure $\mu_t^*$ is approximated from discrete
particle positions.
The optimal transport between discrete measures is not uniquely given by a gradient of
a convex function in the classical sense. The Sinkhorn algorithm
(\Cref{sec:mp-sinkhorn}) bridges this gap: by adding an entropic penalty, the discrete
measures are effectively ``blurred'' into continuous approximations, recovering a
gradient-like structure and enabling stable computation.
\end{remark}

\begin{example}[One-dimensional case: sorted coupling]
\label{ex:mp-1d-wasserstein}
In $d=1$, the optimal transport has a particularly elegant form. For
$\mu, \nu \in \mathcal{P}_2(\mathbb{R})$ with distribution functions $F_\mu, F_\nu$ and
quantile functions $F_\mu^{-1}, F_\nu^{-1}$:
\begin{equation}
    W_2(\mu,\nu)^2 = \int_0^1 |F_\mu^{-1}(u) - F_\nu^{-1}(u)|^2\, du.
    \label{eq:1d-wasserstein}
\end{equation}
The optimal map is $T(x) = F_\nu^{-1}(F_\mu(x))$ --- transport the $u$-quantile of $\mu$
to the $u$-quantile of $\nu$ (sorted coupling). This is computed exactly in
$O(N \log N)$ by sorting the particles. \textbf{The synchronous coupling used in
Chapter 3 for common-noise cancellation can be explicitly constructed in 1-D via this
quantile coupling.}
\end{example}

\subsection{The Sinkhorn Algorithm}
\label{sec:mp-sinkhorn}

For computational purposes, we regularise the Kantorovich problem with an entropic
penalty:
\begin{equation}
    W_{2,\epsilon}(\mu,\nu)^2 := \inf_{\pi \in \Pi(\mu,\nu)} \int |x-y|^2 \pi +
    \epsilon\, \mathrm{KL}(\pi \| \mu \otimes \nu),
    \label{eq:sinkhorn-regularised}
\end{equation}
where $\mathrm{KL}(\pi\|\mu\otimes\nu) = \int \log\frac{d\pi}{d(\mu\otimes\nu)}\, d\pi$.
The solution has the form
\[
    \pi_\epsilon^*(dx,dy) = u(x) K(x,y) v(y)\, \mu(dx)\nu(dy),
\]
where
$K(x,y) = e^{-|x-y|^2/\epsilon}$ is the Gaussian kernel, and $(u,v)$ satisfies the
Sinkhorn fixed-point equations:
\begin{equation}
    u_i = a_i \Big/ \sum_j K_{ij} v_j, \qquad v_j = b_j \Big/ \sum_i K_{ij} u_i,
    \label{eq:sinkhorn-fixed-point}
\end{equation}
where $a_i = \mu(\{x_i\})$, $b_j = \nu(\{y_j\})$. This converges in $O(N^2/\epsilon^2)$
operations (Cuturi, 2013). In Chapter 4, this regularised distance will be used to
compute Wasserstein gradients efficiently.

\subsection{Metric Space Structure}

\begin{theorem}[$(\mathcal{P}_2(\mathbb{R}^d), W_2)$ is a Complete Separable Metric Space]
\label{thm:wasserstein-metric-space}
The Wasserstein-2 distance satisfies:
\begin{enumerate}[label=(\roman*)]
    \item Non-degeneracy: $W_2(\mu,\nu) = 0 \iff \mu = \nu$;
    \item Symmetry: $W_2(\mu,\nu) = W_2(\nu,\mu)$;
    \item Triangle inequality: $W_2(\mu,\rho) \leq W_2(\mu,\nu) + W_2(\nu,\rho)$;
    \item Completeness: every Cauchy sequence in $(\mathcal{P}_2(\mathbb{R}^d), W_2)$
    converges to an element of $\mathcal{P}_2(\mathbb{R}^d)$;
    \item Separability: $\mathcal{P}_2(\mathbb{R}^d)$ contains a countable dense subset.
\end{enumerate}
\end{theorem}

\begin{proof}[Proof of (iii) via gluing lemma]
Given $\mu, \nu, \rho$, let $\pi_{12} \in \Pi(\mu,\nu)$ and $\pi_{23} \in \Pi(\nu,\rho)$
be any couplings. By the gluing lemma (Villani, 2009, Lemma 7.6), there exists
$\pi_{123} \in \mathcal{P}(\mathbb{R}^d \times \mathbb{R}^d \times \mathbb{R}^d)$ with
marginals $(\pi_{12}, \pi_{23})$, i.e., the joint distribution of $(X,Y,Z)$ where
$(X,Y) \sim \pi_{12}$, $(Y,Z) \sim \pi_{23}$. The marginal
$\pi_{13} = (\mathrm{proj}_{13})_\# \pi_{123}$ is a coupling of $\mu$ and $\rho$. By
Minkowski's inequality:
\[
    W_2(\mu,\rho) \leq \left( \int |x-z|^2 \pi_{13}(dx,dz) \right)^{1/2} =
    \mathbb{E}[|X-Z|^2]^{1/2} \leq \mathbb{E}[|X-Y|^2]^{1/2} + \mathbb{E}[|Y-Z|^2]^{1/2}.
\]
Optimising over $\pi_{12}$ and $\pi_{23}$ gives the triangle inequality. The remaining
properties are proved similarly.
\end{proof}

\subsection{Weak Convergence in $\mathcal{P}_2$}

\begin{theorem}[Characterisation of $W_2$-convergence]
\label{thm:w2-characterisation}
The following are equivalent for $\mu_n, \mu \in \mathcal{P}_2(\mathbb{R}^d)$:
\begin{enumerate}[label=(\roman*)]
    \item $W_2(\mu_n, \mu) \to 0$;
    \item $\mu_n \xrightarrow{w} \mu$ weakly and $\int |x|^2\, \mu_n(dx) \to
    \int |x|^2\, \mu(dx)$;
    \item For every continuous $f : \mathbb{R}^d \to \mathbb{R}$ with
    $|f(x)| \leq C(1+|x|^2)$: $\int f\, d\mu_n \to \int f\, d\mu$.
\end{enumerate}
\end{theorem}

Condition (ii) shows that $W_2$-convergence is strictly stronger than weak
convergence: it additionally requires convergence of the second moment. \textbf{This
is precisely the correct topology for the MFG analysis, since $\mu_t = \mathcal{L}(X_t)$
and $\mathbb{E}[|X_t|^2]$ must remain finite and continuous in $t$.}

\subsection{Geodesics and the Lions Derivative}

The space $(\mathcal{P}_2(\mathbb{R}^d), W_2)$ has a rich geometric structure that
motivates the Lions derivative used in the master equation.

\begin{definition}[Wasserstein Geodesic]
The geodesic (shortest path) from $\mu_0$ to $\mu_1$ in $(\mathcal{P}_2(\mathbb{R}^d), W_2)$
is the curve
\[
    \mu_t := \big((1-t)\,\mathrm{id} + t T\big)_\# \mu_0, \qquad t \in [0,1],
\]
where $T = \nabla\varphi$ is the Brenier optimal transport map from $\mu_0$ to $\mu_1$.
This interpolates linearly in position space: the $x$-atom of $\mu_0$ travels in a
straight line to $T(x)$ at constant speed.
\end{definition}

\begin{proposition}[Geodesic is length-minimising]
The path $(\mu_t)_{t \in [0,1]}$ satisfies $W_2(\mu_s, \mu_t) = |t-s|\, W_2(\mu_0,\mu_1)$
for all $s,t \in [0,1]$.
\end{proposition}

The Lions derivative (or $L$-derivative) formalises differentiation of functionals
$F : \mathcal{P}_2(\mathbb{R}^d) \to \mathbb{R}$ along such geodesics. Rather than
being a pointwise derivative, it is a linear functional on the tangent space. For a
functional $F$ that is Lions-differentiable, there exists a measurable function
$D_\mu F(\mu) : \mathbb{R}^d \to \mathbb{R}^d$ (unique up to $\mu$-a.e.\ equality) such
that for any $\nu \in \mathcal{P}_2(\mathbb{R}^d)$ with optimal transport map
$T = \nabla\varphi$ from $\mu$ to $\nu$,
\begin{equation}
    \frac{d}{dt}\Big|_{t=0} F(\mu_t) = \int_{\mathbb{R}^d} D_\mu F(\mu)(x) \cdot (T(x) -
    x)\, \mu(dx).
    \label{eq:lions-derivative}
\end{equation}
In other words, the directional derivative of $F$ along a geodesic with initial
velocity $v = T - \mathrm{id}$ is given by the $L^2(\mu)$ inner product of
$D_\mu F(\mu)$ and $v$. For a comprehensive treatment, see Carmona and Delarue (2018,
Vol.\ I, Chapter 5).

\subsection{Law of Large Numbers in Wasserstein Distance}

\begin{theorem}[LLN Rate in $W_2$, Fournier--Guillin, 2015]
\label{thm:fournier-guillin}
Let $\mu \in \mathcal{P}_2(\mathbb{R}^d)$ and $X_1, X_2, \dots$ be i.i.d.\ with law
$\mu$. Define $\mu^N = \frac{1}{N}\sum_{i=1}^N \delta_{X_i}$.
\begin{enumerate}[label=(\roman*)]
    \item Almost-sure convergence: $W_2(\mu^N, \mu) \to 0$ a.s.\ as $N \to \infty$.
    \item Rate: If $\int_{\mathbb{R}^d} |x|^q\, \mu(dx) < \infty$ for some $q > 4$, then
    for $d=1$:
    \[
        \mathbb{E}[W_2(\mu^N,\mu)] \leq C_\mu \cdot N^{-1/2}.
    \]
    For $d \geq 3$: $\mathbb{E}[W_2(\mu^N,\mu)] \leq C_{\mu,d} \cdot N^{-1/d}$, under
    appropriate moment conditions. (See Fournier and Guillin, 2015, Theorem 1 and
    Corollary 2.)
\end{enumerate}
\end{theorem}

\begin{proof}[Proof of (i) in $d=1$]
By the explicit formula \eqref{eq:1d-wasserstein}:
\[
    W_2(\mu^N,\mu)^2 = \int_0^1 |F_N^{-1}(u) - F^{-1}(u)|^2\, du.
\]
For fixed $u$, $F_N^{-1}(u) \to F^{-1}(u)$ a.s.\ by the Glivenko--Cantelli theorem.
Since $|F_N^{-1}(u) - F^{-1}(u)|^2 \leq 4\max(|F^{-1}(u)|, |F_N^{-1}(u)|)^2$ and the
right-hand side is uniformly integrable by the second moment assumption, the dominated
convergence theorem gives $W_2(\mu^N,\mu)^2 \to 0$ a.s.
\end{proof}

\begin{remark}[Curse of dimensionality and mean-field mitigation]
The rate $N^{-1/d}$ for $d \geq 3$ highlights the curse of dimensionality in
Wasserstein space: for a $d=50$ dimensional state space, the empirical measure converges
exceedingly slowly. \textbf{In practice, the mean-field assumption --- coupled with the
specific convex structure of the rewards and the contraction property established in
Theorem~\ref{thm:wellposedness} --- allows particle approximation (Chapter 5) to circumvent this
theoretical bottleneck.} The Wasserstein distance is used primarily as an analytical
tool for stability, not as the primary object of high-dimensional statistical
estimation.
\end{remark}

\begin{remark}[Significance for the thesis]
\Cref{thm:fournier-guillin} justifies two central approximations: (a) replacing the
$N$-player empirical measure $\mu_t^N$ by the mean-field distribution $\mu_t$ with
error $O(N^{-1/2})$ (for $d=1$ under sufficient moments); (b) validating the
block-bootstrap consistency in Theorem 9.1, where $B = 10{,}000$ samples give
$O(10^{-2.5})$ estimation error.
\end{remark}

\begin{remark}[Wasserstein barycentre -- not used below]
The Wasserstein (Fr\'echet) barycentre of measures $\mu_1,\ldots,\mu_K$ with
weights $\lambda_k$ minimises $\sum_k \lambda_k W_2(\nu,\mu_k)^2$ over $\nu$;
see Agueh and Carlier (2011) for existence/uniqueness and Cuturi and Doucet
(2014) for a Sinkhorn-based algorithm. This concept does not appear elsewhere
in the thesis and is noted here only for completeness.
\end{remark}


\section{Brownian Motion and It\^o Calculus}
\label{sec:mp-ito-calculus}

\subsection{Brownian Motion: Construction and Properties}

\begin{definition}[Standard Brownian Motion]
A standard Brownian motion (Wiener process) $B = (B_t)_{t \geq 0}$ on a probability
space $(\Omega, \mathcal{F}, \mathbb{P})$ is a continuous stochastic process satisfying:
\begin{enumerate}[label=(\roman*)]
    \item $B_0 = 0$ almost surely;
    \item Independent increments: for $0 \leq s < t$, $B_t - B_s$ is independent of
    $\mathcal{F}_s$;
    \item Gaussian increments: $B_t - B_s \sim N(0, t-s)$ for all $0 \leq s < t$;
    \item Continuous paths: $t \mapsto B_t(\omega)$ is continuous for
    $\mathbb{P}$-almost every $\omega$.
\end{enumerate}
\end{definition}

\begin{proposition}[Key path properties]
Brownian motion paths are almost surely:
\begin{enumerate}[label=(\roman*)]
    \item H\"older continuous of order $\alpha < \tfrac{1}{2}$ (but not $\tfrac{1}{2}$);
    \item Nowhere differentiable: $\lim_{h \to 0} \frac{B_{t+h} - B_t}{h}$ does not
    exist for any $t$;
    \item Of unbounded variation: $\sup_\Pi \sum_k |B_{t_k} - B_{t_{k-1}}| = \infty$
    over all partitions $\Pi$;
    \item Of finite quadratic variation: $[B,B]_t = t$ (see below).
\end{enumerate}
\end{proposition}

\subsection{Quadratic Variation}

\begin{definition}[Quadratic Variation]
The quadratic variation of a semimartingale $M$ is the process $[M,M]_t$ defined as
the limit in probability:
\[
    [M,M]_t = \lim_{|\Pi| \to 0} \sum_k (M_{t_k} - M_{t_{k-1}})^2,
\]
where the limit is over partitions $\Pi = \{0 = t_0 < t_1 < \dots < t_n = t\}$ as the
mesh $|\Pi| = \max_k (t_k - t_{k-1}) \to 0$.
\end{definition}

\begin{proposition}[Quadratic variation of Brownian motion]
For standard Brownian motion $B$: $[B,B]_t = t$ a.s.
\end{proposition}

\begin{proof}[Proof sketch]
Let $V_n = \sum_{k=1}^n (B_{kh} - B_{(k-1)h})^2$ with $h = t/n$. Then
$\mathbb{E}[V_n] = n \cdot h = t$ and $\mathrm{Var}(V_n) = n \cdot 2h^2 = 2t^2/n \to 0$.
So $V_n \to t$ in $L^2$, and by Borel--Cantelli, a.s.\ along a subsequence.
\end{proof}

\begin{remark}[Why quadratic variation matters: the It\^o correction]
The fact that $[B,B]_t = t \neq 0$ is the source of all the ``extra terms'' in
stochastic calculus. For a smooth function $f$ and a smooth path $x$,
$f(x_t) = f(x_0) + \int f'(x_s)\, dx_s$ (ordinary chain rule). But for a Brownian path,
the Taylor remainder $\frac{1}{2}f''(B_s)(dB_s)^2$ contributes $\frac{1}{2}f''(B_s)\, ds$
(not zero), because $(dB)^2 = dt$. This is It\^o's formula.
\end{remark}

\subsection{The It\^o Integral}

The classical Riemann--Stieltjes integral $\int_0^T H_t\, dB_t$ is not well-defined
pathwise because $B$ has unbounded variation. It\^o's construction extends it to
square-integrable adapted integrands.

\begin{definition}[It\^o Integral]
Let $(\mathcal{F}_t)$ be the natural filtration of $B$, and $H = (H_t)_{t \in [0,T]}$
be $(\mathcal{F}_t)$-adapted with $\mathbb{E}[\int_0^T H_t^2\, dt] < \infty$. The It\^o
integral $\int_0^T H_t\, dB_t$ is the $L^2(\mathbb{P})$-limit of Riemann sums with left
endpoints:
\[
    \int_0^T H_t\, dB_t := L^2\text{-}\lim_{n \to \infty} \sum_{k=0}^{n-1} H_{t_k}
    (B_{t_{k+1}} - B_{t_k}).
\]
\end{definition}

\begin{theorem}[It\^o Isometry]
\label{thm:ito-isometry}
For adapted square-integrable $H$:
\begin{equation}
    \mathbb{E}\left[ \left( \int_0^T H_t\, dB_t \right)^2 \right] = \mathbb{E}\left[
    \int_0^T H_t^2\, dt \right].
    \label{eq:ito-isometry}
\end{equation}
\end{theorem}

\begin{proof}[Proof sketch]
For simple processes, the expectation of the square expands to
$\sum_k \mathbb{E}[H_{t_k}^2] \Delta t_k$ by independence. The general case follows by
approximation.
\end{proof}

\begin{remark}[Significance for Theorem~\ref{thm:wellposedness}]
The It\^o isometry is used in the proof of Theorem~\ref{thm:wellposedness}: after deriving the Wasserstein
distance ODE $d W_2^2 \leq C W_2^2\, dt + dM_t$, we take expectations and use
$\mathbb{E}[M_t] = 0$ (since $M_t = \int_0^t Z_s\, dW_s^0$ is a martingale by It\^o
isometry, as $\mathbb{E}[\int Z_s^2\, ds] < \infty$). This eliminates the stochastic
term and allows the Gronwall inequality to close.
\end{remark}

\subsection{It\^o's Formula}

\begin{theorem}[It\^o's Formula, one-dimensional]
Let $f \in C^2(\mathbb{R})$ and $X_t = X_0 + \int_0^t b_s\, ds + \int_0^t \sigma_s\, dB_s$
be an It\^o process. Then:
\[
    f(X_t) = f(X_0) + \int_0^t f'(X_s)\, dX_s + \frac{1}{2}\int_0^t f''(X_s)\, d[X,X]_s.
\]
Expanding $d[X,X]_s = \sigma_s^2\, ds$:
\begin{equation}
    f(X_t) = f(X_0) + \int_0^t f'(X_s) b_s\, ds + \int_0^t f'(X_s) \sigma_s\, dB_s +
    \frac{1}{2}\int_0^t f''(X_s) \sigma_s^2\, ds.
    \label{eq:ito-formula-1d}
\end{equation}
\end{theorem}

\begin{theorem}[It\^o's Formula, multidimensional]
\label{thm:ito-multidim}
Let $f \in C^2(\mathbb{R}^d)$ and let $X_t \in \mathbb{R}^d$ satisfy
$dX_t = b_t\, dt + \sigma_t\, dW_t + \sigma_0\, dW_t^0$ (where $W_t \in \mathbb{R}^d$,
$W_t^0 \in \mathbb{R}^{d_0}$, and $\sigma_t \in \mathbb{R}^{d \times d}$,
$\sigma_0 \in \mathbb{R}^{d \times d_0}$). Then:
\begin{align*}
    f(X_t) &= f(X_0) + \int_0^t \nabla f(X_s) \cdot b_s\, ds + \int_0^t \nabla f(X_s)
        \cdot \sigma_s\, dW_s + \int_0^t \nabla f(X_s) \cdot \sigma_0\, dW_s^0 \\
    &\quad + \frac{1}{2}\int_0^t \mathrm{tr}(\sigma_s \sigma_s^\top D^2 f(X_s))\, ds +
        \frac{1}{2}\int_0^t \mathrm{tr}(\sigma_0 \sigma_0^\top D^2 f(X_s))\, ds.
\end{align*}
\end{theorem}

\begin{example}[Applying It\^o to a test function]
\label{ex:mp-ito-test-function}
Let $X_t^i$ be a single agent satisfying $dX_t^i = b_t\, dt + \sigma\, dW_t^i +
\sigma_0\, dW_t^0$. Apply \Cref{thm:ito-multidim} with $f = \phi$ (a smooth test
function):
\[
    \begin{aligned}
    d\phi(X_t^i) &= \nabla\phi(X_t^i) \cdot b_t\, dt + \sigma \nabla\phi(X_t^i) \cdot
    dW_t^i + \sigma_0 \nabla\phi(X_t^i) \cdot dW_t^0 \\
    &\quad + \frac{\sigma^2}{2}
    \Delta\phi(X_t^i)\, dt + \frac{1}{2} \mathrm{tr}(\sigma_0 \sigma_0^\top D^2
    \phi(X_t^i))\, dt.
    \end{aligned}
\]
Averaging over $i = 1, \dots, N$ and writing
$\langle \mu_t^N, \phi \rangle = \frac{1}{N}\sum_i \phi(X_t^i)$:
\[
    \begin{aligned}
    d\langle \mu_t^N, \phi \rangle &= \left\langle \mu_t^N, b_t \cdot \nabla\phi +
    \frac{\sigma^2}{2}\Delta\phi + \frac{1}{2}\mathrm{tr}(\sigma_0 \sigma_0^\top D^2
    \phi) \right\rangle dt \\
    &\quad + \underbrace{\frac{\sigma}{N}\sum_i \nabla\phi(X_t^i)\,
    dW_t^i}_{\to 0 \text{ as } N \to \infty} + \sigma_0 \langle \mu_t^N, \nabla\phi
    \rangle\, dW_t^0.
    \end{aligned}
\]
The idiosyncratic term vanishes in $L^2$ as $N \to \infty$ (mutual independence of
$W^i$). The common noise term survives because $W_t^0$ is the same for all agents.
Taking $N \to \infty$ yields the Fokker--Planck SPDE of \Cref{sec:mp-fokker-planck}.
\end{example}

\subsection{It\^o's Product Rule and Integration by Parts}

\begin{proposition}[It\^o Product Rule]
For It\^o processes $X_t$ and $Y_t$:
\[
    d(X_t Y_t) = X_t\, dY_t + Y_t\, dX_t + d[X,Y]_t,
\]
where $[X,Y]_t = \int_0^t \sigma_s^X \sigma_s^Y\, ds$ is the quadratic covariation
(vanishes if driven by independent Brownians).
\end{proposition}

\begin{example}[Integration by parts for stochastic integrals]
$B_t^2 = 2\int_0^t B_s\, dB_s + t$. (Verify: apply \eqref{eq:ito-formula-1d} with
$f(x) = x^2$, $f'(x) = 2x$, $f''(x) = 2$.)
\end{example}

\subsection{Stochastic Differential Equations}

\begin{theorem}[Existence and Uniqueness for SDEs]
\label{thm:sde-existence-uniqueness}
Let $b : [0,T] \times \mathbb{R}^d \to \mathbb{R}^d$ and
$\sigma : [0,T] \times \mathbb{R}^d \to \mathbb{R}^{d \times m}$ satisfy for some
$L > 0$:
\begin{enumerate}[label=(\roman*)]
    \item Lipschitz: $|b(t,x) - b(t,y)| + \|\sigma(t,x) - \sigma(t,y)\| \leq L|x-y|$;
    \item Linear growth: $|b(t,x)| + \|\sigma(t,x)\| \leq L(1+|x|)$.
\end{enumerate}
Then for any $\mathcal{F}_0$-measurable $X_0$ with $\mathbb{E}[|X_0|^2] < \infty$, the
SDE $dX_t = b(t,X_t)\, dt + \sigma(t,X_t)\, dW_t$ has a unique strong solution
$(X_t)_{t \in [0,T]}$ satisfying
\[
    \mathbb{E}\Big[ \sup_{t \in [0,T]} |X_t|^2 \Big] \leq C(1+\mathbb{E}[|X_0|^2])
    e^{CT}.
\]
\end{theorem}

\begin{proof}[Proof sketch --- Picard iteration]
Define $X_t^0 \equiv X_0$ and
$X_t^{n+1} = X_0 + \int_0^t b(s,X_s^n)\, ds + \int_0^t \sigma(s,X_s^n)\, dW_s$. By
It\^o isometry and the Lipschitz condition:
\[
    \mathbb{E}\Big[ \sup_{s \leq t} |X_s^{n+1} - X_s^n|^2 \Big] \leq C \int_0^t
    \mathbb{E}\Big[ \sup_{r \leq s} |X_r^n - X_r^{n-1}|^2 \Big]\, ds.
\]
Gronwall's inequality gives
$\mathbb{E}[\sup_{s \leq T} |X_s^n - X_s^{n-1}|^2] \leq C^n T^n / n! \to 0$. The Picard
iterates converge in $L^2$ to a unique limit satisfying the SDE.
\end{proof}


\section{Fokker--Planck SPDE}
\label{sec:mp-fokker-planck}

\subsection{The Forward Equation}

The Fokker--Planck equation (also called the Kolmogorov forward equation) governs the
evolution of the law of a diffusion process. In the MFG setting with common noise, it
becomes a stochastic PDE.

\begin{definition}[Fokker--Planck Operator]
For drift $b : \mathbb{R}^d \to \mathbb{R}^d$ and diffusion $\sigma > 0$, the
Fokker--Planck operator is
\[
    \mathcal{L}^* \phi(x) := b(x) \cdot \nabla\phi(x) + \frac{\sigma^2}{2}
    \Delta\phi(x), \qquad \phi \in C_c^\infty(\mathbb{R}^d).
\]
This is the formal $L^2$-adjoint of the generator
$\mathcal{L} f(x) = b(x) \cdot \nabla f(x) + \frac{\sigma^2}{2}\Delta f(x)$:
$\int (\mathcal{L}f)\, d\mu = \int f (\mathcal{L}^* \mu)$ for smooth $\mu$.
\end{definition}

\begin{theorem}[Fokker--Planck SPDE]
\label{thm:fokker-planck-spde}
Let $X_t$ satisfy the McKean--Vlasov SDE: $dX_t = b(t,X_t,\mu_t)\, dt + \sigma\, dW_t^i
+ \sigma_0\, dW_t^0$, $\mu_t = \mathcal{L}(X_t \mid \mathcal{F}_t^{W^0})$. Then $\mu_t$
satisfies the following SPDE in the weak sense: for all $\phi \in C_c^\infty(\mathbb{R}^d)$:
\begin{equation}
    d\langle \mu_t, \phi \rangle = \langle \mu_t, \mathcal{L}_t^* \phi \rangle\, dt +
    \sigma_0 \langle \mu_t, \nabla\phi \rangle\, dW_t^0, \qquad \langle \mu_0, \phi
    \rangle = \int \phi(x)\, \mu_0(dx).
    \label{eq:fp-spde-weak}
\end{equation}
\end{theorem}

\begin{proof}
This follows directly from \Cref{ex:mp-ito-test-function}. By It\^o's formula applied
to $\phi(X_t)$:
\[
    \phi(X_t) - \phi(X_0) = \int_0^t \mathcal{L}_s^* \phi(X_s)\, ds + \sigma \int_0^t
    \nabla\phi(X_s)\, dW_s^i + \sigma_0 \int_0^t \nabla\phi(X_s)\, dW_s^0.
\]
Take the conditional expectation given $\mathcal{F}_t^{W^0}$: the idiosyncratic term
$\sigma \int \nabla\phi(X_s)\, dW_s^i$ has conditional expectation zero (since $W^i$ is
independent of $W^0$). The common noise term $\sigma_0 \int \nabla\phi(X_s)\, dW_s^0$
survives because $W^0$ is measurable with respect to $\mathcal{F}^{W^0}$. This gives
the weak form \eqref{eq:fp-spde-weak}.
\end{proof}

\subsection{Strong Form: The SPDE in Density Form}

When $\mu_t$ has a density $\rho_t(x)$ (i.e., $\mu_t(dx) = \rho_t(x)\, dx$), the SPDE
\eqref{eq:fp-spde-weak} takes the strong form:
\begin{equation}
    d\rho_t = \big[-\nabla_x \cdot (b(t,x,\mu_t)\rho_t) + \frac{\sigma^2}{2}
    \Delta_x \rho_t\big]\, dt + \sigma_0 \nabla_x \rho_t\, dW_t^0.
    \label{eq:fp-spde-strong}
\end{equation}

\begin{remark}[Normalisation and the Kushner--Zakai distinction]
Equation \eqref{eq:fp-spde-strong} is linear in $\rho_t$. This linearity is a
consequence of two facts: (i) the SPDE describes the evolution of the conditional law
of the state given the common noise, and (ii) the total mass is preserved because the
drift and diffusion operators are in divergence form. Consequently, the normalisation
constant remains identically one. This distinguishes the Fokker--Planck SPDE from the
Kushner--Zakai equation of nonlinear filtering, which involves a stochastic
normalisation factor arising from observations. The homogeneous common-noise exposure
(constant $\sigma_0$) is sufficient but not necessary for this linearity; the key
property is that the coefficients do not depend on the measure in a way that would
introduce nonlinear feedback into the normalisation. See Chapter 4 of Carmona and
Delarue (2018) for a detailed discussion.
\end{remark}

\subsection{Sobolev Regularity}

\begin{definition}[Sobolev Spaces]
For $k \in \mathbb{N}$ and $p \in [1,\infty)$:
\[
    W^{k,p}(\mathbb{R}^d) = \{ f \in L^p(\mathbb{R}^d) \mid \partial^\alpha f \in
    L^p(\mathbb{R}^d) \text{ for all } |\alpha| \leq k \},
\]
with norm $\|f\|_{W^{k,p}}^p = \sum_{|\alpha| \leq k} \|\partial^\alpha f\|_{L^p}^p$.
The Hilbert case $p=2$ gives $H^k(\mathbb{R}^d) = W^{k,2}(\mathbb{R}^d)$ with inner
product $\langle f,g \rangle_{H^k} = \sum_{|\alpha| \leq k} \langle \partial^\alpha f,
\partial^\alpha g \rangle_{L^2}$. The dual space is
$H^{-1}(\mathbb{R}^d) = (H^1(\mathbb{R}^d))^*$.
\end{definition}

\begin{theorem}[Fokker--Planck Well-Posedness in Sobolev Space]
\label{thm:fp-sobolev-wellposedness}
Under Assumptions A1--A3 (Lipschitz drift, bounded coefficients, regularity), the SPDE
\eqref{eq:fp-spde-weak} has a unique solution
$\rho \in L^2(\Omega \times [0,T]; H^1(\mathbb{R}^d))$, and
\begin{equation}
    \mathbb{E}\Big[ \sup_{t \in [0,T]} \|\rho_t\|_{L^2}^2 \Big] + \mathbb{E}\left[
    \int_0^T \|\nabla\rho_t\|_{L^2}^2\, dt \right] \leq C(1 + \|\rho_0\|_{H^1}^2).
    \label{eq:fp-sobolev-bound}
\end{equation}
\end{theorem}

\begin{proof}[Proof sketch]
Apply the energy method: multiply \eqref{eq:fp-spde-strong} by $\rho_t$ and integrate
over $\mathbb{R}^d$. Using It\^o's formula for $\|\rho_t\|_{L^2}^2$:
\[
    d\|\rho_t\|_{L^2}^2 = 2\langle \rho_t, d\rho_t \rangle + d[\langle\rho,\cdot\rangle,
    \langle\rho,\cdot\rangle]_t.
\]
The deterministic part contributes $-\sigma^2 \|\nabla\rho_t\|_{L^2}^2\, dt$ (by
integration by parts) plus a term bounded by
$\|b\|_\infty \|\rho_t\|_{L^2} \|\nabla\rho_t\|_{L^2}\, dt$. The stochastic part
contributes $\|\sigma_0\|_F^2 \|\nabla\rho_t\|_{L^2}^2\, dt$ (by It\^o isometry). Taking
expectations and applying Gronwall gives \eqref{eq:fp-sobolev-bound}.
\end{proof}

\subsection{The Gelfand Triple}

\begin{definition}[Gelfand Triple]
A Gelfand triple (or rigged Hilbert space) is a triple of spaces
$V \subset H \subset V^*$, where:
\begin{itemize}
    \item $V$ is a reflexive Banach space (e.g., $H^1(\mathbb{R}^d)$);
    \item $H$ is a Hilbert space (e.g., $L^2(\mathbb{R}^d)$);
    \item $V^*$ is the dual of $V$ (e.g., $H^{-1}(\mathbb{R}^d)$);
    \item the inclusions are continuous and dense.
\end{itemize}
\end{definition}

For the Fokker--Planck SPDE, the relevant Gelfand triple is
$H^1(\mathbb{R}^d) \subset L^2(\mathbb{R}^d) \subset H^{-1}(\mathbb{R}^d)$. The
Fokker--Planck operator $\mathcal{L}^* : H^1(\mathbb{R}^d) \to H^{-1}(\mathbb{R}^d)$ is
bounded with norm $\|\mathcal{L}^*\| \leq L_b + \sigma^2$ (Lemma 4.2 of the thesis,
proved via Hahn--Banach in \Cref{sec:mp-sobolev}). This is what makes the SPDE
well-posed in $H^{-1}$.


\section{Martingale Theory}
\label{sec:mp-martingale}

\subsection{Martingales and Stopping Times}

\begin{definition}[Martingale]
An adapted integrable process $(M_t)_{t \geq 0}$ is a martingale with respect to
$(\mathcal{F}_t)$ if $\mathbb{E}[M_t \mid \mathcal{F}_s] = M_s$ $\mathbb{P}$-a.s.\ for
all $0 \leq s \leq t$. It is a local martingale if there exist stopping times
$\tau_n \uparrow \infty$ such that $(M_{t \wedge \tau_n})$ is a martingale for each $n$.
\end{definition}

\begin{theorem}[It\^o integrals are local martingales]
If $H$ is adapted with $\int_0^t H_s^2\, ds < \infty$ a.s., then
$M_t := \int_0^t H_s\, dB_s$ is a local martingale with $M_0 = 0$. If additionally
$\mathbb{E}[\int_0^T H_s^2\, ds] < \infty$, then $M_t$ is a square-integrable
martingale.
\end{theorem}

\begin{theorem}[Doob's Maximal Inequality]
\label{thm:doob}
For a non-negative submartingale $(M_t)$:
\[
    \mathbb{E}\Big[ \sup_{0 \leq t \leq T} M_t^2 \Big] \leq 4\, \mathbb{E}[M_T^2].
\]
\end{theorem}

\begin{remark}
This is used in Theorem~\ref{thm:wellposedness}: to bound the expected supremum
\[
    \mathbb{E}\big[\sup_t W_2^2(\mu_t^1,\mu_t^2)\big] \leq 4\,\mathbb{E}\big[W_2^2(\mu_T^1,\mu_T^2)\big],
\]
allowing the Gronwall bound to be applied.
\end{remark}

\begin{theorem}[Burkholder--Davis--Gundy (BDG) Inequality]
\label{thm:bdg}
For any continuous local martingale $M$ with $M_0 = 0$ and $p \geq 1$:
\[
    c_p\, \mathbb{E}\big[[M,M]_T^{p/2}\big] \leq \mathbb{E}\Big[ \sup_{t \leq T}
    |M_t|^p \Big] \leq C_p\, \mathbb{E}\big[[M,M]_T^{p/2}\big].
\]
In particular, for $p=2$: $\mathbb{E}[\sup_t M_t^2] \leq 4\mathbb{E}[[M,M]_T]$.
\end{theorem}

\subsection{Martingale Representation Theorem}

\begin{theorem}[Martingale Representation, Brownian Filtration]
\label{thm:martingale-representation}
Let $(\mathcal{F}_t)$ be the natural filtration of a $d$-dimensional Brownian motion
$W$. Then every $(\mathcal{F}_t)$-martingale $M$ has the representation
\[
    M_t = M_0 + \int_0^t H_s\, dW_s
\]
for a unique $(\mathcal{F}_t)$-adapted process $H$ with
$\mathbb{E}[\int_0^T |H_s|^2\, ds] < \infty$.
\end{theorem}

\begin{remark}[Connection to BSDE theory]
The MRT is the existence theorem for solutions to backward SDEs. Given a terminal
condition $\xi = p_T$, the MRT guarantees the existence of $(q_t, r_t)$ such that
\[
    p_T = p_t - \int_t^T \partial_x H_s\, ds + \int_t^T q_s\, dW_s^i + \int_t^T r_s\,
    dW_s^0.
\]
This is the adjoint BSDE from Chapter 1.
\end{remark}

\subsection{The Novikov Condition and Girsanov's Theorem}

\begin{theorem}[Girsanov's Theorem]
\label{thm:girsanov}
Let $\theta : [0,T] \times \Omega \to \mathbb{R}^d$ be adapted with the Novikov
condition
\[
    \mathbb{E}\Big[\exp\big(\tfrac{1}{2}\textstyle\int_0^T |\theta_t|^2\, dt\big)\Big] < \infty.
\]
Define
\[
    Z_T = \exp\left( \int_0^T \theta_t\, dW_t - \frac{1}{2}\int_0^T |\theta_t|^2\, dt
    \right).
\]
Then $Z_T$ is the Radon--Nikodym derivative of a probability measure $\mathbb{P}^*$
with respect to $\mathbb{P}$: $d\mathbb{P}^*/d\mathbb{P} = Z_T$. Under $\mathbb{P}^*$,
the process $\widetilde{W}_t = W_t - \int_0^t \theta_s\, ds$ is a Brownian motion.
\end{theorem}

\begin{remark}[Use in the thesis]
Girsanov's theorem is a standard tool for change of measure in stochastic analysis. In
the MFG setting, it may be applied to remove drift terms under certain monotonicity
assumptions, but the primary analysis in Theorem~\ref{thm:wellposedness} uses synchronous coupling rather
than Girsanov.
\end{remark}


\section{Sobolev Spaces and Functional Analysis}
\label{sec:mp-sobolev}

\subsection{Sobolev Embedding Theorem}

\begin{theorem}[Sobolev Embedding]
\label{thm:sobolev-embedding}
Let $\Omega \subset \mathbb{R}^d$ be bounded with smooth boundary.
\begin{enumerate}[label=(\roman*)]
    \item If $k > d/2$, then $H^k(\Omega) \hookrightarrow C^0(\bar\Omega)$ (continuous
    embedding into continuous functions);
    \item $H^1(\Omega) \hookrightarrow L^{2^*}(\Omega)$ where $2^* = 2d/(d-2)$ for
    $d \geq 3$.
\end{enumerate}
In our setting $\Omega = \mathbb{R}^d$, $d=1$: $H^1(\mathbb{R}) \hookrightarrow
C^{0,1/2}(\mathbb{R})$ (functions are H\"older-$\tfrac{1}{2}$ continuous).
\end{theorem}

\subsection{Key Functional Analysis Tools}
\label{sec:mp-functional-analysis}

Three results from functional analysis are used in the well-posedness theory
(Chapter 3) and the SPDE analysis (Chapter 4). We state each in the form used
and refer to standard references for proofs.

\begin{theorem}[Hahn--Banach Extension, {\citep[Theorem 1.9]{brezis2011}}]
\label{thm:hahn-banach}
Let $E$ be a real normed space, $F \subset E$ a subspace, and $\varphi : F \to
\mathbb{R}$ a bounded linear functional. Then there exists an extension
$\widetilde\varphi : E \to \mathbb{R}$ with $\widetilde\varphi|_F = \varphi$ and
$\|\widetilde\varphi\|_{E^*} = \|\varphi\|_{F^*}$.
\end{theorem}

\noindent\textit{Use here:} shows that the Fokker--Planck operator
$\mathcal{L}^*$, initially defined on smooth test functions (a dense subspace of
$H^1$), extends to a bounded operator
$\mathcal{L}^* : H^1(\mathbb{R}^d) \to H^{-1}(\mathbb{R}^d)$ (Lemma 4.2).

\begin{theorem}[Lumer--Phillips, {\citep[Theorem 3.15]{pazy1983}}]
\label{thm:lumer-phillips}
A closed densely defined operator $A$ on a Banach space $X$ generates a
$C_0$-contraction semigroup if and only if $A$ is dissipative and
$(\mathrm{Id} - A)$ has dense range.
\end{theorem}

\noindent\textit{Use here:} the Fokker--Planck operator $\mathcal{L}^*$ on
$L^2(\mu^*)$ is dissipative ($\langle \mathcal{L}^* \rho, \rho \rangle = -\sigma^2
\|\nabla\rho\|^2 \le 0$), so Lumer--Phillips gives a contraction semigroup. The
spectral gap $\lambda_1 > 0$ (the smallest positive eigenvalue of $-\mathcal{L}^*$)
controls the exponential decay rate of the linearised Fokker--Planck equation.

\begin{theorem}[Log-Sobolev implies Poincar\'e, {\citep[Theorem 5.1]{bakryetal2014}}]
\label{thm:log-sobolev}
If $\mu^*$ satisfies a log-Sobolev inequality with constant $C_{\mathrm{LSI}}$:
$H(\mu \| \mu^*) \le C_{\mathrm{LSI}}\, I(\mu \| \mu^*)$,
then $\lambda_1 \ge 1/C_{\mathrm{LSI}}$.
\end{theorem}

\noindent\textit{Use here:} under the log-Sobolev condition (Assumption A5'),
the spectral gap lower bound $\lambda_1 \ge 1/C_{\mathrm{LSI}}$ gives an
exponential convergence rate for the Fokker--Planck semigroup. As noted in
Chapter 8, this is a strong assumption: log-concave equilibrium measures satisfy
it (Bakry--\'Emery criterion), but heavy-tailed or multimodal equilibria do not.
When only the Poincar\'e inequality holds, the same argument gives algebraic rates
governed by $\lambda_1$.

\section{Stochastic Maximum Principle}
\label{sec:mp-smp}

\subsection{The Control Problem}

We now state the stochastic optimal control problem precisely. Let $A \subset
\mathbb{R}^k$ be the action space (convex compact). A control
$\alpha = (\alpha_t)_{t \in [0,T]}$ is admissible if it is $(\mathcal{F}_t)$-adapted
and $\mathbb{E}[\int_0^T |\alpha_t|^2\, dt] < \infty$. The state $X_t^\alpha \in
\mathbb{R}^d$ satisfies:
\begin{equation}
    dX_t^\alpha = b(t, X_t^\alpha, \mu_t, \alpha_t)\, dt + \sigma\, dW_t^i + \sigma_0\,
    dW_t^0, \qquad X_0^\alpha \sim \mu_0.
    \label{eq:smp-state}
\end{equation}
The objective is to maximise:
\begin{equation}
    J(\alpha) = \mathbb{E}\left[ \int_0^T f(t, X_t^\alpha, \mu_t, \alpha_t)\, dt +
    g(X_T^\alpha, \mu_T) \right].
    \label{eq:smp-objective}
\end{equation}

\subsection{The Hamiltonian}

\begin{definition}[Control Hamiltonian]
The Hamiltonian associated with the control problem is
\begin{equation}
    H(t, x, \mu, p, \alpha) := p \cdot b(t,x,\mu,\alpha) + f(t,x,\mu,\alpha) \in
    \mathbb{R},
    \label{eq:hamiltonian-def}
\end{equation}
where $p \in \mathbb{R}^d$ is the costate variable (adjoint or shadow price). The
maximised Hamiltonian is
\[
    H^*(t,x,\mu,p) := \sup_{\alpha \in A} H(t,x,\mu,p,\alpha).
\]
\end{definition}

Under Assumption A5 (uniform concavity of $H$ in $\alpha$), $H^*$ is attained at a
unique $\alpha^*(t,x,\mu,p)$.

\begin{example}[Linear-quadratic Hamiltonian]
\label{ex:mp-lq-hamiltonian}
For $b = \kappa(\bar{x}-x) + \alpha$ and $f = -x^2 - \tfrac{\lambda}{2}\alpha^2$:
\[
    H(t,x,\mu,p,\alpha) = p(\kappa(\bar{x}-x) + \alpha) - x^2 - \frac{\lambda}{2}
    \alpha^2.
\]
Maximising over $\alpha$: $\partial_\alpha H = p - \lambda\alpha = 0$, so
$\alpha^*(x,\mu,p) = p/\lambda$ and $H^*(t,x,\mu,p) = p\kappa(\bar{x}-x) - x^2 +
\frac{p^2}{2\lambda}$.
\end{example}

\subsection{The Adjoint Process and SMP}

\begin{definition}[Adjoint BSDE]
The adjoint process $(p_t, q_t, r_t)$ associated with an admissible control $\alpha$
and state process $X_t^\alpha$ satisfies the backward SDE:
\begin{equation}
    -dp_t = \partial_x H(t, X_t^\alpha, \mu_t, p_t, \alpha_t)\, dt - q_t\, dW_t^i -
    r_t\, dW_t^0, \qquad p_T = \partial_x g(X_T^\alpha, \mu_T).
    \label{eq:smp-adjoint-bsde}
\end{equation}
Here $p_t \in \mathbb{R}^d$ is the costate, $q_t \in \mathbb{R}^{d \times d}$ accounts
for idiosyncratic noise, and $r_t \in \mathbb{R}^{d \times d_0}$ accounts for common
noise.
\end{definition}

\begin{theorem}[Stochastic Maximum Principle]
\label{thm:smp}
Let $\alpha^*$ be an optimal control for \eqref{eq:smp-objective} and $(p_t,q_t,r_t)$
the associated adjoint process. Then:
\[
    H(t, X_t^*, \mu_t, p_t, \alpha_t^*) = H^*(t, X_t^*, \mu_t, p_t) \qquad
    \mathbb{P}\text{-a.s., a.e.\ } t.
\]
Under Assumption A5 (strict concavity of $H$ in $\alpha$):
\[
    \alpha_t^* = \argmax_{\alpha \in A} H(t, X_t^*, \mu_t, p_t, \alpha) \qquad
    \text{uniquely, } \mathbb{P}\text{-a.s.}
\]
\end{theorem}

\begin{proof}[Proof sketch --- variational method]
Let $\alpha^\epsilon = \alpha^* + \epsilon(\alpha - \alpha^*)$ be a perturbation. The
optimality $J(\alpha^*) \geq J(\alpha^\epsilon)$ gives, after dividing by $\epsilon$ and
taking $\epsilon \to 0$:
\[
    \mathbb{E}\left[ \int_0^T \big(\partial_\alpha H(t, X_t^*, \mu_t, p_t,
    \alpha_t^*)\big) \cdot (\alpha_t - \alpha_t^*)\, dt \right] \geq 0
\]
for all admissible $\alpha$. Since $\alpha$ is arbitrary, $\partial_\alpha H = 0$
$\mathbb{P}$-a.s.\ a.e.\ $t$. Under strict concavity (A5), this uniquely determines
$\alpha_t^*$.
\end{proof}

\subsection{Connection to Hamilton--Jacobi--Bellman}

The SMP and HJB equation are dual formulations of the same problem.

\begin{theorem}[SMP--HJB Connection]
\label{thm:smp-hjb}
If the value function $u(t,x,\mu)$ is $C^{1,2}$ in $(t,x)$ and admits a Lions
derivative in $\mu$, then:
\begin{equation}
    p_t = D_x u(t, X_t^*, \mu_t), \qquad q_t = \sigma D_x^2 u(t, X_t^*, \mu_t), \qquad
    r_t = \sigma_0^\top \mathbb{E}\big[ D_\mu u(t, X_t^*, \mu_t)(\widetilde{X}_t) \mid
    \mathcal{F}_t^{W^0} \big].
    \label{eq:smp-hjb-connection}
\end{equation}
\end{theorem}

Here $D_\mu u(t,x,\mu)(\widetilde{x})$ denotes the Lions derivative of $u$ with respect
to the measure argument, evaluated at the independent copy $\widetilde{X}_t$. This
derivative is defined via the ``lifting'' of functionals on $\mathcal{P}_2(\mathbb{R}^d)$
to functionals on the Hilbert space $L^2(\Omega; \mathbb{R}^d)$; see
Carmona--Delarue (2018, Vol.\ I, Chapter 5) for a comprehensive treatment. The HJB
equation
\[
    -\partial_t u = H^*(t,x,\mu,D_x u) + \bar{b}(t,\mu)\cdot D_\mu u
    + \frac{\sigma^2+\sigma_0^2}{2}\Delta_x u + \sigma_0^2 \,\mathrm{tr}\big(D_x D_\mu
    u\big) + \frac{\sigma_0^2}{2}\int D_\mu^2 u(\widetilde{x},\widetilde{x})\,
    \mu(d\widetilde{x}),
\]
\[
    u(T,x,\mu) = g(x,\mu),
\]
where $\bar{b}(t,\mu)$ is the equilibrium aggregate drift of $\mu_t$ itself (the
population analogue of the individual drift $b$, evaluated at the optimal control).
\textbf{This is the full master equation, including the common-noise diffusion terms
and the Lions-derivative mean-field coupling term $\bar{b}\cdot D_\mu u$; the latter
is frequently omitted in informal expositions but is generally non-zero whenever
$\bar{b}$ depends on $\mu$, as it does in every model considered in this thesis.
Chapter 7's Appendix~A derivation works through this term explicitly for the
linear-quadratic case, including the correction history of an earlier omission;
see the note there for how the coupling term is computed in practice.} The PDE
is the characterisation of the value function, and its unique solution (under
Assumptions A1--A5) gives the adjoint variables via \eqref{eq:smp-hjb-connection}.

\begin{remark}[Significance for the master equation]
Equation \eqref{eq:smp-hjb-connection} reveals why the master equation $U(t,\mu)$
contains derivatives in measure: the adjoint variable $r_t$ involves the Lions
derivative $D_\mu U$, which measures how the value function changes when the measure
$\mu_t$ is perturbed. \textbf{This is the analytical foundation for Chapter 6's
hierarchical RL, where the slow actor learns $D_\mu U$ via the Wonham filter.}
\end{remark}


\section{Summary: Toolkit Reference}
\label{sec:mp-summary}

The following table maps each tool developed in this chapter to its first use in the
thesis.

\begin{table}[ht]
\centering
\small
\caption{Toolkit Reference}
\label{tab:toolkit-reference}
\begin{tabular}{p{2.8cm}p{5cm}p{3.5cm}}
\toprule
\textbf{Tool} & \textbf{Content} & \textbf{First used} \\
\midrule
$\mathcal{P}_2(\mathbb{R}^d)$ & Probability measures with finite 2nd moment & Ch.\ 3
    (state space for $\mu^*$) \\
\addlinespace
Wasserstein distance & $W_2(\mu,\nu) = (\inf_\pi \int |x-y|^2 \pi)^{1/2}$ & Ch.\ 3
    (Theorem~\ref{thm:wellposedness} LHS) \\
\addlinespace
Brenier's Theorem & Optimal transport map $T = \nabla\varphi$ & Ch.\ 3 (coupling in
    Lemma 3.2) \\
\addlinespace
Metric space structure & $(\mathcal{P}_2, W_2)$ complete separable & Ch.\ 3 (Banach
    FPT applies) \\
\addlinespace
$W_2$-convergence characterisation & $W_2 \to 0 \iff$ weak + 2nd moment & Ch.\ 3
    (existence proof) \\
\addlinespace
Geodesic / Lions derivative & Linear functional on $L^2(\mu)$ & Ch.\ 8 (Lions
    derivative) \\
\addlinespace
LLN rate (Fournier--Guillin) & $\mathbb{E}[W_2(\mu^N,\mu)] \leq CN^{-1/2}$ ($d=1$,
    $q>4$) & Ch.\ 7 ($O(1/N)$ proof) \\
\addlinespace
1D Wasserstein formula & $W_2^2 = \int (F_\mu^{-1} - F_\nu^{-1})^2$ & Ch.\ 4 (numerical
    computation) \\
\addlinespace
Brownian motion & Gaussian independent increments & Ch.\ 3 (SDE driving noise) \\
\addlinespace
Quadratic variation & $[B,B]_t = t$ a.s. & It\^o's formula \\
\addlinespace
It\^o isometry & $\mathbb{E}[(\int H dB)^2] = \mathbb{E}[\int H^2 dt]$ & Ch.\ 3 (Step
    4: $\mathbb{E}[M]=0$) \\
\addlinespace
It\^o's formula & $df(X) = f' dX + \tfrac{1}{2}f'' d[X,X]$ & Fokker--Planck derivation \\
\addlinespace
SDE existence/uniqueness & Lipschitz $\Rightarrow$ unique strong solution & Ch.\ 3
    (agent SDE) \\
\addlinespace
It\^o product rule & $d(XY) = X\, dY + Y\, dX + d[X,Y]$ & Ch.\ 3 (Lemma 3.2
    derivation) \\
\addlinespace
Fokker--Planck SPDE & $d\langle\mu_t,\phi\rangle = \langle\mu_t,\mathcal{L}^*\phi\rangle
    dt + \sigma_0\langle\mu_t,\nabla\phi\rangle dW^0$ & Ch.\ 3 (forward equation) \\
\addlinespace
Sobolev well-posedness & $\rho \in L^2(\Omega; H^1)$, energy estimate & Ch.\ 4 (error
    analysis) \\
\addlinespace
Gelfand triple & $H^1 \subset L^2 \subset H^{-1}$ & Ch.\ 4 \\
\addlinespace
Doob's inequality & $\mathbb{E}[\sup M^2] \leq 4\mathbb{E}[M_T^2]$ & Ch.\ 3 (sup
    bound) \\
\addlinespace
BDG inequality & $\mathbb{E}[\sup |M|^p] \sim \mathbb{E}[[M,M]^{p/2}]$ & Ch.\ 3 (moment
    bounds) \\
\addlinespace
Martingale representation & Every martingale = It\^o integral & Ch.\ 3 (BSDE
    existence) \\
\addlinespace
Girsanov's theorem & Change of measure, remove drift & Ch.\ 3 (alternative approach) \\
\addlinespace
Hahn--Banach & Norm-preserving extension & Ch.\ 4 (Lemma 4.2) \\
\addlinespace
Lumer--Phillips & Dissipative $\Rightarrow$ contraction semigroup & Ch.\ 5/6
    (spectral gap) \\
\bottomrule
\end{tabular}
\end{table}
